% SPDX-License-Identifier: GPL-3.0-or-later
% Copyright (C) 2026 Alessandro Greco
\chapter{Ячейки}\label{cap:08-celle}
\aggiornatoal{23 сентября 2026 г.}{коммит torspec \texttt{8e1f245}; tor 0.4.9.13; Arti 2.6.0; протокол соединения 3--5}

\begin{obiettivi}
К концу этой главы вы сможете:
\begin{itemize}[nosep]
  \item описывать побайтно форматы ячеек фиксированной и переменной длины;
  \item объяснять, почему Tor использует единицы фиксированного размера, от чего защищает это решение и от чего не защищает;
  \item интерпретировать CircID и правила их выбора на каждом канале;
  \item распознавать все команды ячеек и основные команды RELAY;
  \item читать заголовок сообщения RELAY и объяснять назначение \texttt{Recognized} и \texttt{Digest};
  \item писать на Python кодировщик/декодировщик ячеек (лабораторная работа~2).
\end{itemize}
\end{obiettivi}

В предыдущей главе мы увидели, как два узла Tor открывают \emph{канал}: TLS-соединение, по которому они согласуют версию протокола и аутентифицируют друг друга. Теперь рассмотрим, \emph{что передаётся} внутри этого канала. Ответ прост: \textbf{ячейки}. Всё, что переносит Tor, от построения цепи до загруженных веб-страниц, разбивается на ячейки. Понимание ячеек даёт нам алфавит, на котором записан весь протокол.

% ---------------------------------------------------------------------------
\section{Зачем нужны ячейки?}\label{sec:celle-perche}

Низколатентная система анонимности не может скрыть сам \emph{факт} общения двух узлов, но может пытаться скрыть \emph{что} они передают и \emph{сколько} промежуточные узлы способны об этом узнать. Если бы каждый узел пересылал данные ровно в том виде, в каком они поступили, пакетами разных размеров, наблюдатель мог бы проследить сообщение длиной 1337 байт через сеть, просто находя сообщение той же длины на каждой линии.

Поэтому исходная архитектура onion routing, которой следует Tor~\cite{dingledine2004tor}, использует \textbf{единицу передачи фиксированного размера}: ячейку. Когда все ячейки имеют одинаковый размер:
\begin{itemize}
  \item длина отдельного сообщения не раскрывает ни его содержимое, ни принадлежность к конкретному потоку;
  \item реле могут просто и эффективно управлять очередями и буферами;
  \item минимальный заголовок позволяет реле переслать ячейку, прочитав всего несколько байтов.
\end{itemize}

\begin{attenzione}[Что ячейки \emph{не} скрывают]
Ячейки стандартизируют \emph{размер} отдельных единиц, но не \emph{количество} единиц и не \emph{время} их передачи. Противник, наблюдающий оба конца цепи, всё ещё может коррелировать объёмы и временные характеристики трафика. Это фундаментальное ограничение Tor, которое мы рассмотрим в \cref{cap:26-attacchi}. Частичные меры защиты --- заполнение (padding) на уровне канала и цепи --- обсуждаются в \cref{cap:07-canali,cap:26-attacchi} и части~V.
\end{attenzione}

\begin{storia}[512 байт]
В статье 2004 года, представившей Tor, описаны ячейки по 512 байт: 2-байтовый идентификатор цепи, 1-байтовая команда и 509 байт данных~\cite{dingledine2004tor}. Тело длиной 509 байт сохранилось до наших дней. Однако идентификатор цепи в версии~4 протокола соединения вырос до 4 байт, увеличив общий размер ячейки до 514 байт.
\end{storia}

% ---------------------------------------------------------------------------
\section{От потока TLS к ячейкам}\label{sec:celle-tls}

После завершения TLS-рукопожатия стороны обмениваются ячейками \emph{последовательно}, одна за другой, как непрерывным потоком байтов~\spec{tor-spec}{Cells (messages on channels)}. Границы TLS-записей и ячеек независимы: одна ячейка может быть разделена между несколькими записями, а одна TLS-запись может содержать несколько ячеек или их фрагментов. Тем не менее спецификация требует: \enquote{разбиение на TLS-записи не должно раскрывать сведений о типе или содержимом ячеек} (перевод).

Для разработчика это имеет практическое следствие: получатель не может ожидать, что сокет будет выдавать \enquote{по одной ячейке}. Нужно накапливать байты в буфере и извлекать ячейку только после получения её целиком, оставляя лишние байты в буфере. Именно так будет работать декодировщик лабораторной работы~2.

% ---------------------------------------------------------------------------
\section{Форматы ячеек}\label{sec:celle-formato}

Существует два формата: \textbf{фиксированной длины} и \textbf{переменной длины}. Оба начинаются одинаковыми полями \texttt{CircID} и \texttt{Command}.

\subsection{Константы}

Спецификация определяет следующие длины~\spec{tor-spec}{Message lengths}, где $v$ --- версия протокола соединения, согласованная на канале:

\begin{center}
\begin{tabular}{lcl}
\toprule
Имя & Байты & Значение \\
\midrule
\file{CELL\_BODY\_LEN} & 509 & тело ячейки фиксированной длины \\
\file{CIRCID\_LEN}$(v)$, $v<4$ & 2 & длина CircID \\
\file{CIRCID\_LEN}$(v)$, $v\ge4$ & 4 & \\
\file{CELL\_LEN}$(v)$, $v<4$ & 512 & размер ячейки фиксированной длины \\
\file{CELL\_LEN}$(v)$, $v\ge4$ & 514 & \\
\bottomrule
\end{tabular}
\end{center}
\noindent Для любого $v$: $\file{CELL\_LEN}(v) = 1 + \file{CIRCID\_LEN}(v) + \file{CELL\_BODY\_LEN}$.

\begin{approfondimento}[Старое имя]
В старом коде и документации \file{CELL\_BODY\_LEN} называется \file{PAYLOAD\_LEN}. Оба имени обозначают одно значение: 509.
\end{approfondimento}

\subsection{Ячейки фиксированной длины}

На \cref{fig:cella-fissa} показана ячейка фиксированной длины с современным протоколом соединения ($v\ge4$).

\begin{figure}[htbp]
\centering
% SPDX-License-Identifier: GPL-3.0-or-later
% Copyright (C) 2026 Alessandro Greco
%% Русское альтернативное описание EPUB.
\altfig{Формат ячейки Tor фиксированной длины для протокола канала v4 и выше: CircID (4 байта) со смещения 0, Command (1 байт) со смещения 4 и Body (509 байт) со смещения 5; следующий за ячейкой байт имеет смещение 514. На схеме условно показаны данные и неиспользованная часть, заполняемая нулями для многих команд с фиксированной длиной. Исключение — RELAY и RELAY\_EARLY: их зашифрованное тело на уровне канала всегда занимает все 509 байт. Ширина полей изображена не в масштабе.}
% Ячейка фиксированной длины, протокол канала >= 4 (ширины полей не в масштабе).
\begin{tikzpicture}[campo/.style={draw,minimum height=1cm,font=\sffamily\small,inner sep=2pt}]
  \node[campo,fill=torviola!15,minimum width=2.6cm] (c) {CircID (4)};
  \node[campo,fill=torviola!35,minimum width=1.6cm,right=0pt of c] (k) {Cmd (1)};
  \node[campo,fill=torverde!15,minimum width=5.2cm,right=0pt of k] (b) {Тело: данные};
  \node[campo,fill=torgrigio!10,minimum width=2.2cm,right=0pt of b] (p) {нули};
  \foreach \n/\o in {c/0,k/4,b/5} \node[font=\scriptsize,above] at (\n.north west) {\o};
  \node[font=\scriptsize,above] at (p.north east) {514};
  \draw[decorate,decoration={brace,mirror,amplitude=5pt}] ([yshift=-2pt]b.south west) -- ([yshift=-2pt]p.south east)
    node[midway,below=6pt,font=\sffamily\small] {\texttt{CELL\_BODY\_LEN} = 509 байт};
\end{tikzpicture}

\caption{Ячейка фиксированной длины (протокол соединения $\ge 4$): 4-байтовый CircID, 1-байтовая команда и 509-байтовое тело, всего 514 байт. Если полезное содержимое короче 509 байт, остаток тела заполняется нулями. Числа обозначают смещения в байтах от начала ячейки.}
\label{fig:cella-fissa}
\end{figure}

Если передаваемая информация занимает меньше 509 байт, отправитель \textbf{обязан} заполнить оставшиеся байты тела нулями, а получатель \textbf{обязан} их игнорировать~\spec{tor-spec}{Padding fixed-length cell bodies}.

\subsection{Ячейки переменной длины}

Некоторым ячейкам, особенно используемым при установлении канала, требуется больше 509 байт, например для сертификата, или значительно меньше. Формат переменной длины решает эту задачу (\cref{fig:cella-variabile}), добавляя 2-байтовое целое поле \texttt{Length} в порядке \emph{big-endian} и не используя заполнение.

\begin{figure}[htbp]
\centering
% SPDX-License-Identifier: GPL-3.0-or-later
% Copyright (C) 2026 Alessandro Greco
%% Русское альтернативное описание EPUB.
\altfig{Ячейка Tor переменной длины после согласования версии протокола канала v4 или выше: CircID занимает 4 байта со смещения 0, Command — 1 байт со смещения 4, Length — 2 байта беззнакового целого в порядке big-endian со смещения 5, Body содержит ровно Length байт со смещения 7. Общий размер — 7 + Length байт, где Length от 0 до 65535; дополнительного заполнения тела нет. Исключение: до согласования VERSIONS всегда использует заголовок v0 с 2-байтовым CircID. Ширины на схеме условны.}
% Ячейка переменной длины после согласования протокола канала v4+ (ширины не в масштабе).
\begin{tikzpicture}[campo/.style={draw,minimum height=1cm,font=\sffamily\small,inner sep=2pt}]
  \node[campo,fill=torviola!15,minimum width=2.6cm] (c) {CircID (4)};
  \node[campo,fill=torviola!35,minimum width=1.6cm,right=0pt of c] (k) {Cmd (1)};
  \node[campo,fill=torviola!55,text=white,minimum width=2cm,right=0pt of k] (l) {Length (2)};
  \node[campo,fill=torverde!15,minimum width=5cm,right=0pt of l] (b) {Тело};
  \foreach \n/\o in {c/0,k/4,l/5,b/7} \node[font=\scriptsize,above] at (\n.north west) {\o};
  \node[font=\scriptsize,above] at (b.north east) {$7+\mathit{Length}$};
  \draw[decorate,decoration={brace,mirror,amplitude=5pt}] ([yshift=-2pt]b.south west) -- ([yshift=-2pt]b.south east)
    node[midway,below=6pt,font=\sffamily\small] {\textit{Length} байт};
\end{tikzpicture}

\caption{Ячейка переменной длины (протокол соединения $\ge 4$): после CircID и команды 2-байтовое big-endian поле Length задаёт число последующих байтов тела. Поэтому тело может содержать от 0 до \num{65535} байт.}
\label{fig:cella-variabile}
\end{figure}

\subsection{Как различать два формата}

Получателю нужно знать, сколько байтов читать, \emph{до} того как он получит всю ячейку. Правило полностью определяется командой~\spec{tor-spec}{Cells (messages on channels)}:
\begin{itemize}
  \item команда 7 (\cella{VERSIONS}) означает ячейку переменной длины;
  \item любая другая команда меньше 128 означает фиксированную длину;
  \item любая команда не меньше 128 означает переменную длину.
\end{itemize}
Иначе говоря, старший бит байта команды служит индикатором формата, за единственным историческим исключением \cella{VERSIONS}.

\begin{storia}[Исключение VERSIONS]
В версии~1 протокола соединения все ячейки имели фиксированную длину. Версия~2 протокола соединения ввела согласование версий, и \cella{VERSIONS}, получившая значение 7 после уже существовавших команд 0--6, стала единственной ячейкой переменной длины. Версия~3 протокола соединения зарезервировала диапазон $\ge 128$ для команд переменной длины. Версии 1 и 2 сегодня устарели, и поддерживать их в реализациях не следует.
\end{storia}

\subsection{Особый случай VERSIONS}

Возникает циклическая зависимость: длина CircID определяется согласованной версией, но сама версия согласуется через ячейку \cella{VERSIONS}. Решение --- \textbf{всегда} отправлять \cella{VERSIONS} так, будто $v=0$, с 2-байтовым CircID. Её тело --- последовательность 2-байтовых big-endian целых, по одному на каждую поддерживаемую версию. Клиент, поддерживающий версии 3, 4 и 5, отправляет:

\begin{center}
\begin{tabular}{lll}
\toprule
Байты & Поле & Значение \\
\midrule
\texttt{00 00} & CircID & 0, 2-байтовый CircID \\
\texttt{07} & Command & \cella{VERSIONS} \\
\texttt{00 06} & Length & тело 6 байт \\
\texttt{00 03 00 04 00 05} & Body & версии 3, 4, 5 \\
\bottomrule
\end{tabular}
\end{center}

Обе стороны выбирают \textbf{наибольшую версию, присутствующую в обеих} ячейках \cella{VERSIONS}. Если общей версии нет или тело содержит нечётное число байтов, соединение должно быть закрыто~\spec{tor-spec}{Negotiating versions with VERSIONS cells}.

\begin{center}
\begin{tabular}{cp{0.8\textwidth}}
\toprule
Версия & Описание \\
\midrule
1 & (устарела) первоначальное установление канала на основе сертификатов \\
2 & (устарела) установление канала с повторным согласованием TLS; вводит ячейки переменной длины \\
3 & (устарела) современное \enquote{внутрипротокольное} установление канала (\cref{cap:07-canali}) \\
4 & 4-байтовый CircID \\
5 & поддержка заполнения на уровне канала и его согласования \\
\bottomrule
\end{tabular}
\end{center}

В современных реализациях tor предлагает версии 3, 4 и 5 (массив \file{or\_protocol\_versions} в \file{src/core/or/connection\_or.c}), а Arti --- только 4 и 5. Поэтому два актуальных экземпляра всегда согласуют версию~5.

% ---------------------------------------------------------------------------
\section{CircID}\label{sec:circid}

Поле \texttt{CircID} идентифицирует цепь, к которой относится ячейка. Его значение равно 0 для ячеек, не принадлежащих ни одной цепи, например \cella{VERSIONS}, \cella{NETINFO} и \cella{PADDING}.

Ключевой момент: \textbf{CircID локальны для канала}. Трёхпереходная цепь пересекает три канала (клиент--guard, guard--middle и middle--exit) и имеет \emph{разный} CircID на каждом. Каждое реле поддерживает таблицу пересылки, сопоставляющую входящую пару (канал, CircID) с исходящей парой (канал, CircID), как показано на \cref{fig:circid-locale}.

\begin{figure}[htbp]
\centering
% SPDX-License-Identifier: GPL-3.0-or-later
% Copyright (C) 2026 Alessandro Greco
%% Русское альтернативное описание EPUB.
\altfig{На схеме цепь клиент — сторожевой — промежуточный — выходной ретранслятор проходит три канала с разными ненулевыми локальными CircID: 0x8000A1C3 на первом, 0x0000712F на втором и 0x93B40E07 на третьем. Таблица сторожевого узла связывает первый идентификатор со вторым, а таблица промежуточного — второй с третьим. На каналах с протоколом v4+ CircID занимает 32 бита; старший бит зависит от того, кто инициировал соединение канала, а не от номера перехода цепи. Сравнение одних CircID не позволяет напрямую сопоставить каналы, но не исключает корреляции по другим признакам трафика.}
% CircID имеет локальную область действия внутри канала.
\begin{tikzpicture}[node distance=2.3cm]
  \node[client] (cl) {Клиент};
  \node[relay,right=of cl] (g) {Сторожевой};
  \node[relay,right=of g] (m) {Промежуточный};
  \node[relay,right=of m] (x) {Выходной};
  \draw[thick,torgrigio] (cl) -- node[above,font=\sffamily\scriptsize]{канал 1} node[below,font=\ttfamily\scriptsize,text=torviola]{0x8000A1C3} (g);
  \draw[thick,torgrigio] (g) -- node[above,font=\sffamily\scriptsize]{канал 2} node[below,font=\ttfamily\scriptsize,text=torviola]{0x0000712F} (m);
  \draw[thick,torgrigio] (m) -- node[above,font=\sffamily\scriptsize]{канал 3} node[below,font=\ttfamily\scriptsize,text=torviola]{0x93B40E07} (x);
  \node[below=0.9cm of g,draw=torgrigio,fill=white,font=\ttfamily\scriptsize,align=left] (tg)
    {(кан.1, 0x8000A1C3)\\\hspace*{1em}$\leftrightarrow$ (кан.2, 0x0000712F)};
  \node[below=0.9cm of m,draw=torgrigio,fill=white,font=\ttfamily\scriptsize,align=left] (tm)
    {(кан.2, 0x0000712F)\\\hspace*{1em}$\leftrightarrow$ (кан.3, 0x93B40E07)};
  \draw[dotted] (g) -- (tg); \draw[dotted] (m) -- (tm);
\end{tikzpicture}

\caption{Одна и та же цепь имеет различный CircID на каждом канале. Guard- и middle-реле переводят идентификатор по таблице пересылки. Поэтому наблюдатель на одной линии не может обнаружить ту же цепь на другой, просто сравнив CircID. Значения иллюстративны, но соблюдают правило старшего бита: на канале~2 бит равен нулю, то есть соединение открыл middle-реле, а guard-реле повторно его использует.}
\label{fig:circid-locale}
\end{figure}

Это решение имеет два преимущества: узлам не нужно координировать выбор идентификаторов, а CircID не становится глобальной меткой, по которой можно связать разные сегменты одной цепи.

\subsection{Кто выбирает CircID?}

CircID выбирает узел, отправляющий ячейку создания цепи (\cella{CREATE2}, \cref{cap:09-circuiti}). Поскольку цепи может создавать любая сторона канала, нужно исключить одновременный выбор одинакового значения. Правило зависит от версии~\spec{tor-spec}{Choosing circuit IDs in create cells}:
\begin{itemize}
  \item \textbf{протокол соединения $\ge 4$}: узел, который \emph{открыл} соединение, использует CircID со старшим битом (MSB), равным 1; другая сторона --- CircID с MSB, равным 0;
  \item \textbf{протокол соединения $\le 3$}: две половины пространства назначались сравнением открытых ключей идентичности двух реле; реле с меньшим ключом использовало MSB 0.
\end{itemize}
Значение 0 зарезервировано; никакое другое значение не зарезервировано. Существующие реализации выбирают CircID \textbf{случайно} из доступных значений, и спецификация рекомендует новым реализациям поступать так же, чтобы не становиться различимыми. Tor прекращает попытки создать новую цепь на канале после 64 последовательных неудач найти свободное значение.

\begin{approfondimento}[Зачем 4 байта?]
При 2 байтах и доступной только половине пространства канал может переносить чуть больше \num{32000} цепей на сторону, а случайный выбор приводит к коллизиям задолго до исчерпания пространства. \proposta{214} (2012) формулирует проблему прямо: \enquote{у реле заканчиваются идентификаторы цепей} (перевод). Предложение ввело версию~4 протокола соединения с 4-байтовыми CircID, реализованную в tor 0.2.4.11-alpha.
\end{approfondimento}

% ---------------------------------------------------------------------------
\section{Команды ячеек}\label{sec:comandi}

Байт \texttt{Command} определяет интерпретацию тела. Следующие таблицы перечисляют все определённые значения~\spec{tor-spec}{Interpreting the fields: Command}. В столбце C \enquote{да} означает, что команде требуется ненулевой CircID, \enquote{нет} --- нулевой. Столбец P показывает первую версию протокола соединения, поддерживающую команду; пустое значение означает как минимум версию~3.

\begin{table}[htbp]
\centering\small
\begin{tabular}{rlccl}
\toprule
Значение & Команда & C & P & Назначение \\
\midrule
0 & \cella{PADDING} & нет & & заполнение на уровне канала \\
1 & \cella{CREATE} & да & & создать цепь (TAP, устарело) \\
2 & \cella{CREATED} & да & & ответ на CREATE (устарело) \\
3 & \cella{RELAY} & да & & end-to-end данные \\
4 & \cella{DESTROY} & да & & уничтожить цепь \\
5 & \cella{CREATE\_FAST} & да & & создать цепь без криптографии с открытым ключом \\
6 & \cella{CREATED\_FAST} & да & & ответ на CREATE\_FAST \\
8 & \cella{NETINFO} & нет & & время и адреса \\
9 & \cella{RELAY\_EARLY} & да & & как RELAY, но в ограниченном количестве \\
10 & \cella{CREATE2} & да & & создать цепь (расширяемый handshake) \\
11 & \cella{CREATED2} & да & & ответ на CREATE2 \\
12 & \cella{PADDING\_NEGOTIATE} & да & 5 & согласование заполнения \\
\bottomrule
\end{tabular}
\caption{Команды ячеек фиксированной длины.}
\label{tab:comandi-fissi}
\end{table}

\begin{table}[htbp]
\centering\small
\begin{tabular}{rlcl}
\toprule
Значение & Команда & C & Назначение \\
\midrule
7 & \cella{VERSIONS} & нет & согласование протокола соединения \\
128 & \cella{VPADDING} & нет & заполнение переменной длины \\
129 & \cella{CERTS} & нет & сертификаты \\
130 & \cella{AUTH\_CHALLENGE} & нет & вызов для аутентификации \\
131 & \cella{AUTHENTICATE} & нет & аутентификация инициатора \\
132 & \cella{AUTHORIZE} & нет & зарезервировано \\
\bottomrule
\end{tabular}
\caption{Команды ячеек переменной длины.}
\label{tab:comandi-variabili}
\end{table}

\paragraph{Неизвестные команды.} Другие значения не разрешены. Получатель неизвестной команды \emph{может} молча отбросить ячейку и \emph{может} закрыть канал. Такая строгость связана с тем, что каждая новая команда вводится с новой версией протокола соединения: корректная реализация никогда не отправит команду, поддержку которой другая сторона не объявляла. До появления версионирования протокола соединения неизвестные команды, напротив, игнорировались в предположении, что они пришли от более новой версии Tor.

Команды можно сгруппировать по назначению:
\begin{description}
  \item[Управление каналом] \cella{VERSIONS}, \cella{CERTS}, \cella{AUTH\_CHALLENGE}, \cella{AUTHENTICATE}, \cella{NETINFO}: установление канала из \cref{cap:07-canali}.
  \item[Заполнение] \cella{PADDING}, \cella{VPADDING}, \cella{PADDING\_NEGOTIATE}: фиктивный трафик, затрудняющий анализ трафика на линии.
  \item[Жизненный цикл цепи] \cella{CREATE2}/\cella{CREATED2} и их исторические варианты, \cella{DESTROY}: создание и уничтожение цепей (\cref{cap:09-circuiti}).
  \item[Транспорт] \cella{RELAY} и \cella{RELAY\_EARLY}: всё, что передаётся \emph{внутри} цепи.
\end{description}

\subsection{DESTROY}

Тело ячейки \cella{DESTROY} содержит всего один байт, задающий \emph{причину} закрытия цепи~\spec{tor-spec}{Tearing down circuits}. Получив \cella{DESTROY}, узел освобождает ресурсы цепи и, если сам не является конечной точкой, передаёт ячейку \cella{DESTROY} следующему узлу.

\begin{center}\small
\begin{tabular}{rl@{\qquad}rl}
\toprule
0 & NONE & 7 & OR\_IDENTITY \\
1 & PROTOCOL & 8 & CHANNEL\_CLOSED \\
2 & INTERNAL & 9 & FINISHED \\
3 & REQUESTED & 10 & TIMEOUT \\
4 & HIBERNATING & 11 & DESTROYED \\
5 & RESOURCELIMIT & 12 & NOSUCHSERVICE \\
6 & CONNECTFAILED & & \\
\bottomrule
\end{tabular}
\end{center}

\begin{attenzione}[Код причины как побочный канал]
Спецификация рекомендует \textbf{всегда} отправлять причину NONE и не передавать далее полученный код причины. Если реле пересылало бы исходную причину в обоих направлениях, два сговорившихся реле одной цепи могли бы использовать её как сигнал, чтобы определить, что находятся в одной цепи: маленький пример того, как даже один байт способен стать \emph{побочным каналом}.
\end{attenzione}

\subsection{RELAY\_EARLY}

\cella{RELAY\_EARLY} имеет тот же формат, что \cella{RELAY}, но служит для \textbf{ограничения длины цепи}~\spec{tor-spec}{Handling RELAY\_EARLY cells}. Правила:
\begin{itemize}
  \item клиент должен отправлять сообщения расширения цепи (\texttt{EXTEND2}) только внутри \cella{RELAY\_EARLY}; начиная с версии 0.2.3.11-alpha реле отклоняют расширения, полученные иначе;
  \item реле, получившее более 8 прямых \cella{RELAY\_EARLY} в одной цепи, должно, как правило, закрыть её; при получении хотя бы одной \emph{обратной}, направленной к клиенту, ячейки оно обязано немедленно закрыть цепь;
  \item чтобы не раскрывать длину цепи, клиенты также должны, как правило, отправлять примерно первые 8 ячеек, не предназначенных первому переходу, как \cella{RELAY\_EARLY}.
\end{itemize}
Без этого ограничения противник мог бы строить чрезвычайно длинные цепи, многократно проходящие через одни и те же реле и тем самым усиливающие его собственный трафик. Это атака перегрузки длинными маршрутами, описанная Эвансом, Динглдайном и Гротхоффом~\cite{evans2009congestion}.

\begin{storia}[Неожиданный сигнал]
Правило, запрещающее обратные \cella{RELAY\_EARLY}, появилось после инцидента 2014 года: некоторые реле отправляли такие ячейки клиентам как скрытый канал для маркировки цепей пользователей onion-сервисов. Подробно этот случай разбирается в \cref{cap:27-casi-rete}.
\end{storia}

% ---------------------------------------------------------------------------
\section{Внутри RELAY-ячеек}\label{sec:relay}

Ячейки \cella{RELAY} и \cella{RELAY\_EARLY} несут \textbf{сообщения RELAY}: команды, которыми обмениваются клиент и узлы цепи, от расширения цепи до данных TCP-потока. Их 509-байтовые тела зашифрованы слоями, по одному на каждый переход. Это и есть \enquote{луковица} из \cref{cap:10-cifratura-relay}. С точки зрения канала тело RELAY-ячейки поэтому выглядит полностью заполненным: нулевого padding не видно.

После снятия всех слоёв шифрования открытое тело следует традиционному формату на \cref{fig:relay}~\spec{tor-spec}{Relay cells}.

\begin{figure}[htbp]
\centering
% SPDX-License-Identifier: GPL-3.0-or-later
% Copyright (C) 2026 Alessandro Greco
%% Русское альтернативное описание EPUB.
\altfig{Расшифрованное тело традиционной ячейки RELAY (формат tor1): заголовок 11 байт содержит Cmd (1 байт, смещение 0), Recognized (2, смещение 1), StreamID (2, смещение 3), Digest (4, смещение 5) и Length (2, смещение 9). Данные занимают Length байт от смещения 11, не более 498 байт; остаток до смещения 509 при необходимости заполняется. При отправке значение Recognized равно нулю; традиционный скользящий дайджест включает и заполнение. Спецификация рекомендует сначала до четырёх нулевых байтов, а затем случайные байты. Если осталось менее четырёх байтов, все они нулевые; при Length = 498 заполнения нет. Получатель не должен полагаться на его содержимое.}
% Расшифрованное тело RELAY-ячейки, традиционный формат.
\begin{tikzpicture}[campo/.style={draw,minimum height=1cm,font=\sffamily\scriptsize,inner sep=1pt,align=center}]
  \node[campo,fill=torviola!35,minimum width=1.2cm] (a) {Cmd\\(1)};
  \node[campo,fill=torviola!15,minimum width=1.6cm,right=0pt of a] (b) {Recognized\\(2)};
  \node[campo,fill=torviola!15,minimum width=1.4cm,right=0pt of b] (c) {StreamID\\(2)};
  \node[campo,fill=torviola!15,minimum width=1.4cm,right=0pt of c] (d) {Digest\\(4)};
  \node[campo,fill=torviola!15,minimum width=1.3cm,right=0pt of d] (e) {Length\\(2)};
  \node[campo,fill=torverde!15,minimum width=3.2cm,right=0pt of e] (f) {Data\\(Length)};
  \node[campo,fill=torgrigio!10,minimum width=2cm,right=0pt of f] (g) {Padding\\нули + случайные};
  \foreach \n/\o in {a/0,b/1,c/3,d/5,e/9,f/11} \node[font=\scriptsize,above] at (\n.north west) {\o};
  \node[font=\scriptsize,above] at (g.north east) {509};
  \draw[decorate,decoration={brace,mirror,amplitude=5pt}] ([yshift=-2pt]a.south west) -- ([yshift=-2pt]e.south east)
    node[midway,below=6pt,font=\sffamily\small] {заголовок: 11 байт};
  \draw[decorate,decoration={brace,mirror,amplitude=5pt}] ([yshift=-2pt]f.south west) -- ([yshift=-2pt]g.south east)
    node[midway,below=6pt,font=\sffamily\small] {498 байт};
\end{tikzpicture}

\caption{\emph{Расшифрованное} тело традиционной RELAY-ячейки: заголовок 11 байт, до 498 байт данных и, при наличии места, заполнение до 509 байт. Спецификация рекомендует до четырёх нулевых байтов вначале, затем случайные байты; если осталось менее четырёх байтов, все они нулевые, а при 498 байтах данных заполнения нет.}
\label{fig:relay}
\end{figure}

\begin{description}
  \item[Команда RELAY (1 байт)] Тип сообщения (\cref{tab:comandi-relay}).
  \item[Recognized (2 байта)] Ноль в открытом сообщении. Это поле позволяет реле быстро определить, адресована ли ячейка ему: если после снятия собственного слоя поле ненулевое, ячейка наверняка всё ещё зашифрована и должна быть переслана.
  \item[StreamID (2 байта)] Поток, к которому относится сообщение, выбранный клиентом; ноль для сообщений о цепи в целом.
  \item[Digest (4 байта)] Первые 4 байта \emph{накопительного} дайджеста всех сообщений, которыми обменивались с этим переходом (SHA-1 для обычных цепей, SHA3-256 для цепей onion-сервисов). Он подтверждает, что ячейка предназначена этому реле и не была изменена.
  \item[Length (2 байта)] Число значимых байтов в поле Data.
  \item[Data] Содержимое сообщения.
  \item[Заполнение] Остаток, доводящий тело до 509 байт.
\end{description}

11-байтовый заголовок оставляет максимум $509 - 11 = 498$ байт данных на ячейку.

\subsection{Recognized и Digest: двухступенчатый фильтр}

Зачем нужны два поля? Вычислять дайджест для каждой транзитной ячейки было бы дорого. Recognized --- дешёвый фильтр: после расшифрования ячейки, которая всё ещё зашифрована, его 2 байта псевдослучайны и равны нулю лишь с вероятностью $2^{-16}$. Только при нулевом Recognized реле вычисляет дайджест. Если дайджест совпадает, ячейка \enquote{распознана} и обрабатывается; иначе пересылается дальше. Вычисление и его слабости рассматриваются в \cref{cap:10-cifratura-relay}.

\subsection{Заполнение сообщений RELAY}

Заполнение сообщений RELAY отличается от заполнения нулями ячеек фиксированной длины. Спецификация рекомендует 4 нулевых байта, а затем \textbf{случайные байты}~\spec{tor-spec}{Relay cells}. Случайные байты делают содержимое ячейки непредсказуемым. Это необходимо для аутентифицированных сообщений SENDME из \proposta[Предложения]{289}, которые доказывают фактическое получение данных, включая дайджест полученных ячеек (\cref{cap:11-stream}). Получатель не должен зависеть от содержимого заполнения.

\subsection{Команды RELAY}

\Cref{tab:comandi-relay} перечисляет определённые команды RELAY. Столбец Dir. задаёт направление: F (\emph{forward}, только от клиента внутрь цепи), B (\emph{backward}, только к клиенту) или F/B (оба направления). C означает \emph{управляющее} сообщение, для которого StreamID должен быть нулевым.

\begin{table}[htbp]
\centering\small
\begin{tabularx}{\textwidth}{@{}rll>{\raggedright\arraybackslash}X@{}}
\toprule
Значение & Команда & Dir. & Назначение \\
\midrule
\multicolumn{4}{l}{\emph{Основной протокол}} \\
1 & BEGIN & F & открыть поток \\
2 & DATA & F/B & данные потока \\
3 & END & F/B & закрыть поток \\
4 & CONNECTED & B & поток успешно открыт \\
5 & SENDME & F/B, C? & подтверждение получения (управление потоком) \\
6--7 & EXTEND/EXTENDED & F/B, C & расширение с TAP (устарело) \\
8--9 & TRUNCATE/TRUNCATED & F/B, C & усечение цепи (не используется) \\
10 & DROP & F/B, C & заполнение между концами цепи \\
11--12 & RESOLVE/RESOLVED & F/B & удалённое разрешение DNS \\
13 & BEGIN\_DIR & F & поток к кэшу каталогов \\
14--15 & EXTEND2/EXTENDED2 & F/B, C & расширение цепи \\
16--18 & \emph{зарезервировано} & & UDP поверх Tor (\proposta{339}) \\
\multicolumn{4}{l}{\emph{Conflux}} \\
19--22 & CONFLUX\_\ldots & F/B, C & LINK, LINKED, LINKED\_ACK, SWITCH: несколько цепей (\proposta{329}) \\
\multicolumn{4}{l}{\emph{Onion-сервисы}} \\
32 & ESTABLISH\_INTRO & F, C & создать точку введения \\
33 & ESTABLISH\_RENDEZVOUS & F, C & создать точку встречи \\
34--35 & INTRODUCE1/INTRODUCE2 & F/B, C & запрос введения \\
36--37 & RENDEZVOUS1/RENDEZVOUS2 & F/B, C & соединить стороны в rendezvous \\
38--40 & \ldots\_ESTABLISHED, INTRODUCE\_ACK & B, C & подтверждения \\
\multicolumn{4}{l}{\emph{Заполнение цепи и управление потоком}} \\
41--42 & PADDING\_NEGOTIATE/\ldots D & F/B, C & заполнение цепи \\
43--44 & XOFF/XON & F/B & управление потоком (\proposta{324}) \\
\bottomrule
\end{tabularx}
\caption{Команды RELAY. SENDME является управляющим сообщением только тогда, когда относится к цепи, а не к отдельному потоку (\cref{cap:11-stream}).}
\label{tab:comandi-relay}
\end{table}

Если ячейка распознана, но её relay command неизвестна, её необходимо отбросить; тем не менее её содержимое учитывается в дайджесте и окнах управления потоком.

% ---------------------------------------------------------------------------
\section{К новому формату: предложение 340 и CGO}\label{sec:nuovo-formato}

Традиционный формат имеет два известных ограничения. Первое связано с эффективностью: большое сообщение всегда занимает как минимум целую ячейку, тогда как несколько маленьких сообщений нельзя упаковать в одну. Второе связано с безопасностью: 4-байтовый дайджест и неаутентифицированное шифрование AES-CTR оставляют возможность атак маркирования (\emph{tagging}), которые мы рассмотрим в \cref{cap:10-cifratura-relay}.

\proposta{340} (\emph{Packed and fragmented relay messages}) разделяет понятия \emph{ячейки} и \emph{сообщения}: одно сообщение может фрагментироваться по нескольким ячейкам, а несколько сообщений --- упаковываться в одну. \proposta{359} (\emph{Counter Galois Onion, Updated}) вводит новую схему шифрования RELAY, CGO, и использует формат ячейки предложения 340. Расшифрованная CGO-ячейка имеет структуру:

\begin{center}
\begin{tabular}{ll}
\toprule
Поле & Байты \\
\midrule
тег аутентификации & 16 \\
тело & $509 - 16 = 493$ \\
\midrule
\quad \emph{внутри тела:} команда & 1 \\
\quad длина & 2 \\
\quad StreamID (только для некоторых команд) & 2 \\
\quad данные & \emph{длина} \\
\quad заполнение (4 нуля, затем случайные байты) & остаток \\
\bottomrule
\end{tabular}
\end{center}

StreamID присутствует только для команд, которым он требуется (BEGIN, DATA, END, CONNECTED, RESOLVE, RESOLVED, BEGIN\_DIR, XON, XOFF). Recognized и Digest исчезают, их заменяет тег аутентификации. Клиент использует CGO с реле только тогда, когда его поддерживают обе стороны (субпротокол \texttt{Relay=6}) и согласуют его при создании цепи.

\begin{approfondimento}[Сколько данных помещается в CGO-ячейку?]
Из формата следует, что для сообщения DATA доступно $493 - 1 - 2 - 2 = 488$ байт данных, на 10 меньше традиционного формата: это цена 16-байтового тега аутентификации.
\end{approfondimento}

CGO уже не просто предложение: tor поддерживает его начиная с серии 0.4.9 (введён в 0.4.9.3-alpha, сентябрь 2025 года). В Arti поддержка всегда включена; параметр сборки \texttt{counter-galois-onion} сохраняется лишь для совместимости. Репозиторий спецификаций по-прежнему помечает оба предложения как \enquote{Open}: статус предложения отражает процесс стандартизации и не всегда совпадает с состоянием кода.

% ---------------------------------------------------------------------------
\section{Анатомия сеанса}\label{sec:sessione}

Теперь можно собрать части вместе и ячейка за ячейкой проследить, что клиент отправляет guard-реле, чтобы построить трёхпереходную цепь и загрузить веб-страницу. Последовательность показана на \cref{fig:sessione}; в следующих главах каждый шаг разбирается подробно.

\begin{figure}[htbp]
\centering
% SPDX-License-Identifier: GPL-3.0-or-later
% Copyright (C) 2026 Alessandro Greco
%% Русское альтернативное описание EPUB.
\altfig{Схема обмена на TLS-канале клиент — сторожевой узел. После TLS обе стороны отправляют VERSIONS с CircID 0; одна стрелка сторожевого узла объединяет четыре отдельные ячейки VERSIONS, CERTS, AUTH\_CHALLENGE и NETINFO. Клиент отправляет NETINFO, создаёт первый переход с CREATE2/CREATED2 и расширяет цепь до промежуточного и выходного узлов двумя RELAY\_EARLY [EXTEND2] и ответами RELAY [EXTENDED2]. BEGIN показан в RELAY\_EARLY лишь как допустимый пример, не обязательное правило; затем следуют CONNECTED, запрос и ответ DATA, при необходимости SENDME и завершение DESTROY. Сторожевой узел наблюдает ячейки своего канала, но не читает сообщения для следующих узлов. Названия в скобках поясняют логические команды.}
% Последовательность ячеек на канале клиент — сторожевой узел.
\begin{tikzpicture}[font=\sffamily\small,>=Stealth]
  \node[client] (c) at (0,0) {Клиент};
  \node[relay] (g) at (7,0) {Сторожевой};
  \draw[torgrigio] (c.south) -- ++(0,-12.2);
  \draw[torgrigio] (g.south) -- ++(0,-12.2);
  \newcommand{\av}[3]{\draw[cellaf] (0,#1) -- node[above,font=\ttfamily\footnotesize]{#2} (7,#1);
    \node[font=\scriptsize,text=torgrigio,anchor=west] at (7.4,#1) {#3};}
  \newcommand{\ind}[3]{\draw[risposta] (7,#1) -- node[above,font=\ttfamily\footnotesize]{#2} (0,#1);
    \node[font=\scriptsize,text=torgrigio,anchor=west] at (7.4,#1) {#3};}
  \av{-1.1}{VERSIONS}{CircID 0}
  \ind{-1.8}{VERSIONS, CERTS, AUTH\_CHALLENGE, NETINFO}{CircID 0}
  \av{-2.5}{NETINFO}{канал готов}
  \av{-3.5}{CREATE2}{CircID 0x8\ldots, 1-й переход}
  \ind{-4.2}{CREATED2}{}
  \av{-5.2}{RELAY\_EARLY [EXTEND2]}{к промежуточному}
  \ind{-5.9}{RELAY [EXTENDED2]}{}
  \av{-6.9}{RELAY\_EARLY [EXTEND2]}{к выходному}
  \ind{-7.6}{RELAY [EXTENDED2]}{}
  \av{-8.6}{RELAY\_EARLY [BEGIN]}{поток к сайту}
  \ind{-9.3}{RELAY [CONNECTED]}{}
  \av{-10.1}{RELAY [DATA]}{HTTP-запрос}
  \ind{-10.8}{RELAY [DATA] \ldots}{ответ, SENDME}
  \av{-11.7}{DESTROY}{цепь завершается}
\end{tikzpicture}

\caption{Ячейки, которыми обмениваются клиент и guard-реле по TLS-каналу при создании трёхпереходной цепи и открытии потока. Guard-реле видит только ячейки собственного канала: оно не может читать содержимое RELAY-ячеек, предназначенных middle- и exit-реле, поскольку они зашифрованы ключами, которыми guard-реле не владеет.}
\label{fig:sessione}
\end{figure}

\begin{enumerate}
  \item \textbf{Установление канала}: \cella{VERSIONS}, затем \cella{CERTS}, \cella{AUTH\_CHALLENGE} и \cella{NETINFO} от guard-реле; клиент отвечает \cella{NETINFO}. CircID всегда равен 0.
  \item \textbf{Первый переход}: \cella{CREATE2} с новым CircID (MSB 1, поскольку соединение открыл клиент); guard-реле отвечает \cella{CREATED2}.
  \item \textbf{Второй и третий переходы}: две ячейки \cella{RELAY\_EARLY}, содержащие \texttt{EXTEND2}; guard-реле, а затем middle-реле преобразуют их в ячейки \cella{CREATE2}, отправляемые следующему узлу. Ответы возвращаются как \cella{RELAY} с \texttt{EXTENDED2}.
  \item \textbf{Поток}: ячейка с \texttt{BEGIN} для exit-реле, которое отвечает \texttt{CONNECTED}; затем \texttt{DATA} в обоих направлениях вперемешку с \texttt{SENDME}.
  \item \textbf{Закрытие}: \texttt{END} для потока, а в конце жизни цепи --- \cella{DESTROY}.
\end{enumerate}

% ---------------------------------------------------------------------------
\section{Лабораторная работа: кодировщик ячеек}\label{sec:lab-celle}

\begin{laboratorio}[2]{Кодирование и декодирование ячеек}
\textbf{Цель.} Написать модуль \file{labs/ru/код/torlab/cells.py}, который будет использоваться во всех последующих лабораторных работах.\\
\textbf{Предварительные условия.} Python~3.10 или новее, \cmd{pytest}. Сеть не нужна: лабораторная работа оперирует вручную построенными байтами.\\
\textbf{Оценочное время.} 1--2 часа.
\end{laboratorio}

\paragraph{Шаг 1: константы и команды.} Превратим таблицы спецификации в перечисления Python, чтобы каждое числовое значение задавалось только в одном месте:

\begin{lstlisting}[language=Python]
CELL_BODY_LEN = 509
RELAY_HEADER_LEN = 11
RELAY_DATA_MAX = CELL_BODY_LEN - RELAY_HEADER_LEN   # 498

class CellCommand(IntEnum):
    PADDING = 0
    CREATE = 1
    # ... все значения из таблиц 8.1 и 8.2
    AUTHORIZE = 132

def circid_len(version):
    return 2 if version < 4 else 4

def is_variable_length(command):
    return command == CellCommand.VERSIONS or command >= 128
\end{lstlisting}

\paragraph{Шаг 2: кодирование.} Функция получает CircID, команду, тело и версию link protocol и выдаёт байты для отправки. Обратите внимание на особую обработку \cella{VERSIONS} и нулевой padding:

\begin{lstlisting}[language=Python]
def encode(cell, version):
    if cell.command == CellCommand.VERSIONS:
        version = 0                  # CircID всегда 2 байта
    n = circid_len(version)
    header = cell.circ_id.to_bytes(n, "big") + bytes([cell.command])
    if is_variable_length(cell.command):
        return (header + struct.pack(">H", len(cell.body))
                + cell.body)
    return header + cell.body.ljust(CELL_BODY_LEN, b"\x00")
\end{lstlisting}

\paragraph{Шаг 3: декодирование потока.} Поскольку байты приходят из TLS-сокета фрагментами, декодировщик возвращает \texttt{None}, пока полной ячейки нет. Иначе он возвращает ячейку и неиспользованные байты:

\begin{lstlisting}[language=Python]
def decode(buffer, version):
    n = circid_len(version)
    if len(buffer) < n + 1:
        return None
    circ_id = int.from_bytes(buffer[:n], "big")
    command = buffer[n]
    if is_variable_length(command):
        if len(buffer) < n + 3:
            return None
        (length,) = struct.unpack(">H", buffer[n+1:n+3])
        start, end = n + 3, n + 3 + length
    else:
        start, end = n + 1, n + 1 + CELL_BODY_LEN
    if len(buffer) < end:
        return None
    return Cell(circ_id, command, buffer[start:end]), buffer[end:]
\end{lstlisting}

\paragraph{Шаг 4: VERSIONS и согласование версии.}

\begin{lstlisting}[language=Python]
def negotiate_version(ours, received_body):
    if len(received_body) % 2:
        raise CellError("VERSIONS содержит нечётное число байтов")
    theirs = {v for (v,) in struct.iter_unpack(">H", received_body)}
    common = set(ours) & theirs
    if not common:
        raise CellError("нет общей версии link protocol")
    return max(common)
\end{lstlisting}

Диагностические строки в русской лабораторной среде локализованы, поскольку предназначены читателю; числовые значения протокола и API функции остаются неизменными.

\paragraph{Шаг 5: сообщение RELAY.} Наконец, строим 11-байтовый заголовок. Строка формата \texttt{">BHH4sH"} для \texttt{struct} точно соответствует полям на \cref{fig:relay}:

\begin{lstlisting}[language=Python]
def encode_relay(msg, random_bytes=os.urandom):
    header = struct.pack(">BHH4sH", msg.command, msg.recognized,
                               msg.stream_id, msg.digest, len(msg.data))
    space = CELL_BODY_LEN - RELAY_HEADER_LEN - len(msg.data)
    padding = b"\x00" * min(4, space) + random_bytes(max(0, space - 4))
    return header + msg.data + padding
\end{lstlisting}

Параметр \texttt{random\_bytes} позволяет тестам заменять случайные байты предсказуемыми значениями.

\paragraph{Проверка.} Репозиторий содержит тесты в \file{labs/ru/код/torlab/test/test\_cells.py}. Запустите их из каталога \file{labs/ru/код}:
\begin{lstlisting}[language=bash,numbers=none]
$ python3 -m pytest -q torlab/test/test_cells.py
........                                                  [100%]
8 passed
\end{lstlisting}
Без имени файла \cmd{pytest} запускает тесты всех лабораторных работ репозитория.
Например, один тест проверяет побайтное совпадение ячейки \cella{VERSIONS} для версий 3, 4 и 5 с примером из \cref{sec:celle-formato}: \texttt{0000 07 0006 0003 0004 0005}.

\paragraph{Расширения.} В лабораторной работе~1 (\cref{cap:07-canali}) этот модуль используется для общения с настоящим реле частной сети. В лабораторной работе~4 мы добавим вычисление digest и послойное шифрование.

% ---------------------------------------------------------------------------
\section*{Итоги}
\begin{itemize}
  \item Ячейки --- единицы передачи Tor по каналу. Начиная с версии~4 протокола соединения ячейка фиксированной длины содержит 514 байт (CircID 4 + команда 1 + тело 509); старые версии использовали 512 байт.
  \item Ячейки переменной длины добавляют 2-байтовое поле Length. \cella{VERSIONS} (7) и все команды $\ge 128$ имеют переменную длину.
  \item \cella{VERSIONS} всегда использует 2-байтовый CircID; обе стороны выбирают наибольшую общую версию.
  \item CircID локален для каждого канала. Сторона, открывшая соединение, выбирает случайные значения с MSB 1.
  \item \cella{DESTROY} должна, как правило, всегда нести причину NONE, чтобы не создавать побочные каналы; \cella{RELAY\_EARLY} ограничивает длину цепи.
  \item Расшифрованное тело RELAY-ячейки имеет 11-байтовый заголовок (command, Recognized, StreamID, Digest, Length) и максимум 498 байт данных. Padding состоит из 4 нулевых, затем случайных байтов.
  \item Предложения 340 и 359 (CGO) вводят новый формат с 16-байтовым тегом аутентификации и необязательным StreamID.
\end{itemize}

% ---------------------------------------------------------------------------
\section*{Упражнения}

\begin{esercizio}[1]\label{es:08-lunghezze}
Сколько байт занимает ячейка \cella{CREATE2} на канале, согласовавшем link protocol~3? А версию~5? Сколько байт занимает \cella{VERSIONS} с телом \texttt{00 04 00 05}?
\end{esercizio}

\begin{esercizio}[1]\label{es:08-formato}
Для каждого значения команды 0, 7, 9, 12, 129 и 200 укажите, имеет ячейка фиксированную или переменную длину. Какое из этих значений недопустимо?
\end{esercizio}

\begin{esercizio}[1]\label{es:08-circid}
Клиент открывает канал к guard-реле с link protocol~5. Какой из CircID он может использовать для новой цепи: \texttt{0x00000000}, \texttt{0x12345678} или \texttt{0x92345678}? Почему?
\end{esercizio}

\begin{esercizio}[2]\label{es:08-hex}
Вручную декодируйте следующую последовательность байтов, полученную на канале с link protocol~4. Для каждой ячейки укажите CircID, команду и длину тела:\\
\texttt{00 00 00 00 81 00 03 aa bb cc 80 00 00 07 04 00 \ldots}\\
(Вторая ячейка продолжается оставшейся частью тела, которая не показана.)
\end{esercizio}

\begin{esercizio}[2]\label{es:08-relay}
Запишите в шестнадцатеричном виде первые 11 байт открытого тела relay message DATA в потоке 5, содержащего 300 байт данных, до вычисления digest. Сколько байт padding следует после данных и сколько из них гарантированно равны нулю?
\end{esercizio}

\begin{esercizio}[2]\label{es:08-recognized}
Промежуточное реле обрабатывает \num{1000000} ячеек, которые ему не адресованы. Для скольких из них в среднем потребуется вычислять digest? Почему одного поля Recognized недостаточно, чтобы определить, предназначена ли ячейка этому реле?
\end{esercizio}

\begin{esercizio}[2]\label{es:08-destroy}
Опишите конкретный сценарий, в котором два сговорившихся реле могли бы использовать код причины в \cella{DESTROY}, чтобы узнать, что находятся в одной цепи. Почему использование причины NONE предотвращает этот сценарий?
\end{esercizio}

\begin{esercizio}[2]\label{es:08-efficienza}
Клиент загружает через цепь файл размером \SI{1}{\mebi\byte} (\num{1048576} байт). Сколько ячеек DATA требуется в традиционном формате? Сколько с CGO? В каждом случае какой процент полезных данных приходится на накладные расходы заголовков ячеек (link protocol~4, без учёта TLS)?
\end{esercizio}

\begin{esercizio}[3]\label{es:08-lab-fuzz}
Расширьте лабораторную работу~2 \emph{property-based} тестом, например с библиотекой Hypothesis, который генерирует случайные ячейки фиксированной и переменной длины, кодирует и объединяет их, а затем делит результат на фрагменты случайной длины. Проверьте, что декодирование возвращает в точности исходные ячейки.
\end{esercizio}

\begin{esercizio}[3]\label{es:08-arti}
Найдите в исходном коде Arti функции кодирования и декодирования ячеек канала (crate \texttt{tor-cell}). Сравните обработку \enquote{неполной ячейки} с функцией \texttt{decode} лабораторной работы~2.
\end{esercizio}

% ---------------------------------------------------------------------------
\section*{Дополнительная литература}
\begin{itemize}
  \item \titolo{Tor Protocol Specification}, разделы \enquote{Cells}, \enquote{Relay cells} и \enquote{Negotiating channels}~\cite{torspec}: первичный источник для этой главы.
  \item R.~Dingledine, N.~Mathewson, P.~Syverson, \titolo{Tor: The Second-Generation Onion Router}~\cite{dingledine2004tor}, раздел 4: исходный формат ячеек и его мотивация.
  \item N.~Evans, R.~Dingledine, C.~Grothoff, \titolo{A Practical Congestion Attack on Tor Using Long Paths}~\cite{evans2009congestion}: атака, мотивировавшая \cella{RELAY\_EARLY}.
  \item Предложения 340 и 359 в torspec: будущий формат ячеек.
\end{itemize}
